\begin{abstract}
FragPicker (SOSP '21) defragments only the file regions an application's requests
straddle, using an analysis phase that turns a BCC trace into per-file lists of hot
regions. A later fork of its artifact added a parallel analysis mode and reported a
12.74$\times$ speedup with 16 threads, measured with a 5\,ms \texttt{sleep} added in
user space before every per-file extent lookup. We re-evaluate that claim with real
traces, an ablation ladder from the unmodified code to the parallel mode, and delays
injected at the block layer instead of in user space. The parallel mode parallelises
an extent lookup that the original analysis never performs and whose result it
discards; on file systems without FIEMAP it drops every file. The original's cost is
process creation (fourteen new processes per traced file) and one full tree walk per
file. A 200-line single-threaded rewrite with byte-identical output is
\textbf{[X]}$\times$ faster than the original on real traces and beats every parallel
configuration. With the delay at the block layer, where only reads that miss the cache
are slowed, the extent lookups of files created together need \textbf{[R]} device reads
in all, and threads gain \textbf{[S]}$\times$ instead of \textbf{[Z]}$\times$. We also
find that the shipped tracing procedure loses the last buffered chunk of events and
wraps file positions above 4\,GiB. All scripts, traces and raw data are released.
\end{abstract}
